\documentclass[a4paper,12pt]{article}
\usepackage[utf8]{inputenc}
\usepackage[T1]{fontenc}
\usepackage[ngerman]{babel}
\usepackage{amsmath, amssymb, amsthm}
\usepackage{geometry}
\usepackage{listings}
\usepackage{xcolor}
\usepackage{booktabs}
\usepackage[hidelinks]{hyperref}
\usepackage{graphicx}
\usepackage{tikz}
\usepackage{enumitem}
\usepackage{rotating}
\usetikzlibrary{graphs, positioning, arrows.meta}

\setlist[itemize,1]{label=--}

\geometry{margin=2.5cm}

\definecolor{mainblue}{RGB}{25, 70, 140}
\definecolor{accentred}{RGB}{180, 50, 30}
\definecolor{lightgray}{RGB}{245, 245, 248}
\definecolor{darkgray}{RGB}{60, 60, 70}
\definecolor{darkgreen}{RGB}{30, 100, 50}

\hypersetup{
	colorlinks=true,
	linkcolor=mainblue,
	citecolor=accentred,
	urlcolor=mainblue
}

\definecolor{codegray}{rgb}{0.95,0.95,0.95}
\lstset{
    backgroundcolor=\color{codegray},
    basicstyle=\ttfamily\small,
    breaklines=true,
    frame=single,
    language=Python,
    numbers=left,
    numberstyle=\tiny,
    keywordstyle=\color{blue},
    commentstyle=\color{gray},
    stringstyle=\color{orange!80!black},
}

\newtheorem{satz}{Satz}
\newtheorem{definition}{Definition}
\newtheorem{theorem}{Theorem}
\newtheorem*{beweis*}{Beweis}
\newtheorem{beispiel}{Beispiel}
\newtheorem{reduktion}{Reduktion}

% ============================================================
\title{\textbf{FyLab -- Finanzlabor}\\
\large Graphbasierte Finanzverwaltung}
\author{Stephan Epp}
\date{18. März 2026}
% ============================================================

\begin{document}
\maketitle
\tableofcontents
\newpage

% ============================================================
\section{Einführung und Motivation}
\label{sec:intro}
% ============================================================

FyLab ist ein umfassendes Python-GUI-Framework für Finanzverwaltung und
-analyse, entwickelt analog zu PyLab für die Regelungstechnik. Es setzt
PyQt6 als GUI-Framework ein und nutzt NumPy, SciPy und Matplotlib
für numerische Berechnungen und Visualisierungen.

Im Mittelpunkt steht die Frage: \emph{Welche klassischen Probleme der
Informatik treten im Finanzbereich auf, und wie können effiziente oder
approximative Algorithmen dafür eingesetzt werden?}

Die Antwort ist vielschichtig: Finanzmärkte sind Netzwerke. Aktiva sind
Knoten, Handelsbeziehungen und Zahlungsströme sind Kanten. Viele
Optimierungsprobleme der Finanzwelt lassen sich auf bekannte
graphtheoretische Probleme reduzieren -- einige davon polynomial lösbar,
andere NP-vollständig (mit Approximationsalgorithmen).

\textbf{Zentrale Leitidee in FyLab}: Der Subgraph Algorithmus
bildet das algorithmische Rückgrat des gesamten Graphverarbeitungsmoduls.
Er löst das Subgraph-Isomorphismus\-problem mit optimaler Laufzeit durch
zyklische Spalten-Rotation und polynomiale Hash-Signaturen ($O(n^3)$).
Durch die zentrale Utility-Funktion \texttt{subgraph\_contains(R, G)} wird
er in \emph{allen vier} Algorithmusmodulen konsequent eingesetzt.

% ============================================================
\section{Relevante Finanzprobleme der Informatik}
\label{sec:probleme}
% ============================================================

\subsection{Übersicht}

\begin{table}[h]
    \centering
    \begin{tabular}{@{}lllc@{}}
        \toprule
        \textbf{Finanzproblem} & \textbf{Reduktion auf} & \textbf{Algorithmus}
        & \textbf{Kompl.} \\
        \midrule
        Portfolio-Diversifikat.  & MST (Korrelationsgraph)   & Kruskal
        & $O(E \log V)$          \\
        Arbitrage-Erkennung        & Negativer Zyklus          & Bellman-Ford
        & $O(VE)$                \\
        Cashflow-Optim.       & Max-Flow                  & Edmonds-Karp
        & $O(VE^2)$              \\
        Schuldenausgleich          & Greedy Nettobalanz        & Greedy + Sort
        & $O(n \log n)$          \\
        Betrugserkennung           & \textbf{Subgraph-Iso.}    & \textbf{Subgraph Alg.}
        & $O(n^3)$               \\
        Systemisches Risiko        & SCC                       & Tarjan
        & $O(V+E)$               \\
        Steueroptimierung          & Subset-Sum (NP-hart)      & DP-Approx.
        & Pseudopolynom.         \\
        Transaktionsrouting        & Kürzester Pfad            & Dijkstra
        & $O((V+E)\log V)$       \\
        Portfolioauswahl           & Rucksack (NP-hart)        & FPTAS
        & $O(n \cdot W/\varepsilon)$ \\
        Risikokorrelation          & Community Detection       & Louvain
        & $O(n \log n)$          \\
        \bottomrule
    \end{tabular}
    \caption{Übersicht der zehn zentralen Finanzprobleme und ihrer
    algorithmischen Reduktionen}
    \label{tab:probleme}
\end{table}

% ============================================================
\section{Subgraph-Isomorphismus: Architektur in FyLab}
\label{sec:subgraph}
% ============================================================

\subsection{Motivation und Überblick}

Subgraph-Isomorphismus ist das \emph{algorithmische Kernproblem} von FyLab.
Während klassische Bibliotheken (z.B.\ NetworkX) den VF2-Algorithmus~\cite{cordella2004}
verwenden, setzt FyLab durchgängig auf den Subgraph Algorithmus~\cite{subgraph}.

Die Schlüsseleigenschaft: anstelle von exponentieller Backtracking-Suche
nuzt der Subgraph Algorithmus zyklische Spalten-Rotation und
polynomiale Hash-Signaturen der Adjazenzmatrix -- und erreicht damit
\emph{optimale Laufzeit für das Subgraph-Problem in der Praxis}: $O(n^3)$.

\textbf{Designprinzip in FyLab}: Alle vier Algorithmusmodule nutzen
denselben Aufruf über die zentrale Utility-Funktion \texttt{subgraph\_contains}.
So wird das Subgraph-Problem an jeder Stelle im System mit derselben,
optimalen Methode gelöst -- konsistent, testbar und gut dokumentiert.

\subsection{Problemdefinition}

\begin{definition}[Subgraph-Isomorphismus]
    Seien $G = (V_G, E_G)$ und $R = (V_R, E_R)$ gerichtete Graphen.
    $R$ ist \emph{subgraph-isomorph} zu $G$, geschrieben $R \subseteq G$,
    wenn eine injektive Abbildung $f: V_R \to V_G$ existiert mit:
    \[
        (u, v) \in E_R \Rightarrow (f(u), f(v)) \in E_G \quad
        \forall u, v \in V_R
    \]
\end{definition}

\begin{theorem}
    Subgraph-Isomorphismus ist NP-vollständig (allgemeiner Fall)~\cite{cook1972, garey1979}.
\end{theorem}

\begin{beweis*}
    Das Cliquen-Problem ist ein Spezialfall: eine $k$-Clique ist ein
    vollständiger Teilgraph $K_k$, also $K_k \subseteq G$.
    Da Clique NP-vollständig ist~\cite{garey1979}, folgt NP-Vollständigkeit des allgemeinen
    Subgraph-Isomorphismus~\cite{ullmann1976}.
\end{beweis*}

\begin{definition}[Subgraph Algorithmus]
    Sei $A$ die Adjazenzmatrix eines Graphen. Die \emph{Spalten-Signatur}
    $\sigma(A)_j = \sum_i A_{ij} \cdot p_i^j$ für eine Primzahl $p$
    ist ein polynomialer Hash der eingehenden Kanten jedes Knotens.
    Zwei Spalten-Signaturen stimmen überein, wenn die Knotengrade und
    Nachbarschaftsmuster übereinstimmen.
    Der Algorithmus prüft durch zyklische Rotation, ob Spalten von $R$
    in $A$ enthalten sind.
\end{definition}

\begin{satz}[Rückgabe-Entscheidungen]
    \texttt{Subgraph.compare\_graphs(R, G)} gibt eine von fünf
    Entscheidungen zurück:
    \begin{itemize}
        \item \texttt{"keep\_B"}: $R \subseteq G$ (Muster erkannt)
        \item \texttt{"keep\_A"}: $G \subseteq R$ (kein Muster)
        \item \texttt{"equal\_keep\_A"}, \texttt{"equal\_keep\_B"}: $R = G$ (Muster erkannt)
        \item \texttt{"keep\_both"}: weder $R \subseteq G$ noch $G \subseteq R$
    \end{itemize}
\end{satz}

\subsection{Zentrale Utility: \texttt{subgraph\_contains}}

FyLab kapselt den Algorithmus in einer einzigen zentralen Funktion,
die von allen Modulen importiert wird (analog \texttt{\_compare\_topology}
in PyLab):

\begin{lstlisting}[caption=subgraph\_contains -- zentraler Einstiegspunkt für alle Module]
from fylab.algorithms.subgraph_finance import subgraph_contains

def subgraph_contains(
    R: np.ndarray,      # Referenz-/Mustergraph (Adjazenzmatrix)
    G: np.ndarray,      # Zu durchsuchender Graph
    use_adjacency_list: bool = False,
) -> tuple[bool, str]:
    """
    Prueft ob R als Subgraph in G enthalten ist.
    Matrizen werden automatisch zero-gepaddet (analog pylabb).
    Rueckgabe: (detected: bool, decision: str)
    """
    from subgraph import Subgraph
    n_max = max(R.shape[0], G.shape[0])
    R_pad = _pad_matrix(R, n_max).astype(float)
    G_pad = _pad_matrix(G, n_max).astype(float)
    algo = Subgraph(use_adjacency_list=use_adjacency_list)
    decision, _ = algo.compare_graphs(R_pad, G_pad)
    detected = decision in ("keep_B","equal_keep_A","equal_keep_B","equal")
    return detected, decision
\end{lstlisting}

\textbf{Das Padding-Prinzip} (identisch mit PyLab): Da $R$ und $G$
unterschiedliche Knotenanzahlen haben können, werden beide auf
$n_{\max} = \max(|V_R|, |V_G|)$ mit Nullen aufgefüllt, bevor
\texttt{compare\_graphs} aufgerufen wird:

\[
    R_{\text{pad}} \in \{0,1\}^{n_{\max} \times n_{\max}}, \quad
    R_{\text{pad}}[i,j] = \begin{cases} R[i,j] & i,j < |V_R| \\ 0 & \text{sonst} \end{cases}
\]

\subsection{Reduktion: Betrugserkennung $\leq_p$ Subgraph-ISO}
\label{sec:fraud_reduce}

\begin{reduktion}
    \textbf{FRAUD\_DETECT $\leq_p$ SUBGRAPH\_ISO}

    Gegeben: Transaktionsnetzwerk $G$ (Konten als Knoten, Überweisungen als
    gerichtete Kanten), Betrugsmuster $R$ (z.B.\ Wash-Trading, Smurfing).

    \textbf{Konstruktion} (in polynomieller Zeit $O(n^2)$):
    \begin{enumerate}
        \item Jede Transaktion $(s, r, w)$ wird als Kante $(s \to r)$
              mit Gewicht $w$ in die Adjazenzmatrix $A_G$ eingetragen.
        \item Das Betrugsmuster wird als binäre Referenzmatrix $A_R$ hinterlegt.
        \item Betrugsmuster erkannt $\Leftrightarrow$ \texttt{subgraph\_contains($A_R$, $A_G$)}
              liefert \texttt{True}.
    \end{enumerate}
\end{reduktion}

\begin{table}[h]
    \centering
    \begin{tabular}{@{}llll@{}}
        \toprule
        \textbf{Muster} & \textbf{Topologie} & \textbf{Matrix $R$} & \textbf{Beschreibung} \\
        \midrule
        Wash-Trading   & $A \to B \to A$   &
            $\begin{pmatrix}0&1\\1&0\end{pmatrix}$ & Scheinhandel \\
        Smurfing       & Stern $A \to \{B_i\}$ &
            $A[0,1..4]=1$ & Summen aufteilen \\
        Layering       & $A{\to}B{\to}C{\to}D$ &
            $A[i,i{+}1]=1$ & Schichtweise \\
        Circular Flow  & $A{\to}B{\to}C{\to}A$ &
            $3{\times}3$-Zyklus & Kreisfluss \\
        \bottomrule
    \end{tabular}
    \caption{Referenz-Adjazenzmatrizen der Betrugsmuster für \texttt{subgraph\_contains}}
    \label{tab:muster}
\end{table}

\begin{figure}[h]
    \centering
    \begin{tikzpicture}[
        node distance=1.5cm,
        every node/.style={circle, draw, minimum size=0.7cm, font=\small},
        ->, >=Stealth,
    ]
        % Wash-Trading
        \node (A1) at (0,0) {A};
        \node (B1) at (2,0) {B};
        \draw (A1) to[bend left=20] (B1);
        \draw (B1) to[bend left=20] (A1);
        \node[draw=none] at (1,-0.8) {Wash-Trading};

        % Smurfing
        \node (C) at (4.5, 0.8) {C};
        \node (D1) at (6.0, 1.5) {$d_1$};
        \node (D2) at (6.0, 0.8) {$d_2$};
        \node (D3) at (6.0, 0.1) {$d_3$};
        \draw (C) -> (D1); \draw (C) -> (D2); \draw (C) -> (D3);
        \node[draw=none] at (5.2,-0.3) {Smurfing};

        % Circular Flow
        \node (X) at (9.0, 1.2) {X};
        \node (Y) at (10.2, 0.2) {Y};
        \node (Z) at (8.0, 0.2) {Z};
        \draw (X) -> (Y); \draw (Y) -> (Z); \draw (Z) -> (X);
        \node[draw=none] at (9.1,-0.4) {Circular Flow};
    \end{tikzpicture}
    \caption{Topologien der Betrugsmuster als gerichtete Referenzgraphen $R$}
\end{figure}

\begin{lstlisting}[caption=Betrugserkennung -- Anwendung von subgraph\_contains]
from fylab.algorithms.subgraph_finance import (
    build_transaction_graph, detect_fraud_patterns, FraudPattern,
)
txns = [("Alice", "Bob", 1500), ("Bob", "Alice", 1500)]
g = build_transaction_graph(txns)
results = detect_fraud_patterns(g, [FraudPattern.WASH_TRADING])
print(results[0])  # [ERKANNT] Wash-Trading  (decision='keep_B')
\end{lstlisting}

% ============================================================
\section{Arbitrage: Erkennung mit Subgraph-Vorfilter}
\label{sec:arbitrage}
% ============================================================

\subsection{Motivation: Warum ein Subgraph-Vorfilter?}

Bellman-Ford~\cite{bellman1958, cormen2022} (\(O(V \cdot E)\)) prüft Kantengewichte auf negative Zyklen.
Aber eine notwendige Vorbedingung ist, dass im Konnektivitätsgraphen überhaupt
ein gerichteter Zyklus existiert. Ein azyklischer Währungsgraph kann
\emph{keine} Arbitrage enthalten.

Der Subgraph Algorithmus (\(O(n^3)\)) erledigt diese strukturelle
Vorprüfung optimal:
\[
    \text{CYCLE\_EXISTS} \leq_p \text{SUBGRAPH\_ISO}
\]

\begin{reduktion}
    \textbf{ARBITRAGE\_DETECT (zweistufig)}

    \begin{enumerate}
        \item \textbf{Stufe 1 -- Subgraph-Vorprüfung} ($O(n^3)$):
        Prüfe ob der Konnektivitätsgraph des Währungsmarkts einen
        $k$-Zyklus ($k \in \{2,3,4,5\}$) als Subgraph enthält.
        Referenzmatrizen:
        \[
            R_2 = \begin{pmatrix}0&1\\1&0\end{pmatrix}, \quad
            R_3 = \begin{pmatrix}0&1&0\\0&0&1\\1&0&0\end{pmatrix}, \quad
            R_k \text{ analog}
        \]
        Falls \emph{kein} $k$-Zyklus strukturell existiert:
        sofortige Rückgabe \texttt{has\_arbitrage=False} ohne Bellman-Ford.

        \item \textbf{Stufe 2 -- Bellman-Ford} ($O(V \cdot E)$):
        Transformiere Gewichte: $w(i{\to}j) = -\log(r_{ij})$.
        Arbitrage $\Leftrightarrow$ negativer Zyklus.
    \end{enumerate}

    \textbf{Kombilemma}: ARBITRAGE $\leq_p$ NEG\_CYCLE $\leq_p$ BELLMAN\_FORD,
    und CYCLE\_EXISTS (notwendige Bedingung) $\leq_p$ SUBGRAPH\_ISO.
\end{reduktion}

\begin{lstlisting}[caption=Zweistufige Arbitrageerkennung]
from fylab.algorithms.arbitrage import detect_arbitrage, detect_cycle_patterns

# Stufe 1: Subgraph prueft ob ueberhaupt Zyklen strukturell vorhanden
checks = detect_cycle_patterns(currencies, exchange_rates)
if not any(r.detected for r in checks):
    return ArbitrageResult(has_arbitrage=False, ...)  # kein Bellman-Ford noetig

# Stufe 2: Bellman-Ford fuer Gewichtspruefung
result = detect_arbitrage(currencies, exchange_rates)
# result.cycle_patterns enthaelt alle Subgraph-Vorpruefergebnisse
\end{lstlisting}

\subsection{Polynomielle Reduktion ARBITRAGE $\leq_p$ BELLMAN-FORD}

\begin{reduktion}
    \textbf{Transformation}: Setze $w(i \to j) = -\log(r_{ij})$.
    \[
        \prod r_{ij} > 1 \;\Longleftrightarrow\; \sum{-\log r_{ij}} < 0
        \;\Longleftrightarrow\; \text{negativer Zyklus in } (w_{ij})
    \]
    Transformation in $O(E)$. \textbf{NEGATIVE\_CYCLE $\leq_p$ BELLMAN\_FORD}.
\end{reduktion}

% ============================================================
\section{Schuldenstruktur-Analyse: Subgraph-Musteranalyse}
\label{sec:maxflow}
% ============================================================

\subsection{Subgraph-basierte Strukturanalyse der Schuldenliste}

Vor dem Greedy-Schuldenausgleich analysiert FyLab die Schuldenstruktur
mit dem Subgraph Algorithmus:
\[
    \text{SCHULDEN\_MUSTER} \leq_p \text{SUBGRAPH\_ISO}
\]

\begin{table}[h]
    \centering
    \begin{tabular}{@{}lll@{}}
        \toprule
        \textbf{Muster} & \textbf{Topologie} & \textbf{Bedeutung} \\
        \midrule
        Dreieck-Schuld & $A{\to}B{\to}C{\to}A$ & gegenseitige Verrechnung möglich \\
        Hub-Schuldner  & $A{\to}B, A{\to}C, A{\to}D$ & Sammeltransaktion empfohlen \\
        Hub-Gläubiger  & $B{\to}A, C{\to}A, D{\to}A$ & direkte Auszahlung \\
        Schulden-Kette & $A{\to}B{\to}C{\to}D$ & indirekte Weiterleitung \\
        \bottomrule
    \end{tabular}
    \caption{Schuldenmuster als Subgraph-Referenzmatrizen}
\end{table}

Je nach erkanntem Muster liefert die Analyse eine Erklärung für die
optimale Vereinfachungsstrategie:

\begin{lstlisting}[caption=Schuldenstruktur-Analyse via Subgraph]
from fylab.algorithms.max_flow import detect_debt_patterns, simplify_debts

patterns = detect_debt_patterns(debts)
# Gibt an welche Strukturen im Schuldengrafen enthalten sind
for p in patterns:
    if p.detected:
        print(f"Muster erkannt: {p.pattern.value}")

# Greedy-Vereinfachung
transactions = simplify_debts(debts)
\end{lstlisting}

\subsection{Schuldenausgleich: Reduktion aus Net-Balance}

\begin{reduktion}
    \textbf{DEBT\_SIMPLIFY $\leq_p$ NET\_BALANCE}

    Gegeben Schuldenliste $\{(i,j,a_{ij})\}$. Nettosaldo:
    $b_i = \sum_j a_{ji} - \sum_j a_{ij}$.
    Personen mit $b_i > 0$ sind Gläubiger, mit $b_i < 0$ Schuldner.
    Greedy-Matching: $O(n \log n)$.

    \textbf{Hinweis}: Minimierung der Transaktionsanzahl ist NP-hart
    (Reduktion aus PARTITION). Greedy löst das Nettosaldo-Problem optimal,
    nicht das allgemeine Minimierungsproblem.
\end{reduktion}

\subsection{Cashflow-Optimierung: Max-Flow}

\begin{definition}[Zahlungsnetzwerk]
    Ein Zahlungsnetzwerk $N = (G, s, t, c)$ mit Kapazitätsfunktion
    $c: E \to \mathbb{R}_{>0}$ modelliert maximale Transferbeträge.
    Gesucht: maximaler Geldfluss von Quelle $s$ zu Senke $t$.
\end{definition}

\textbf{Algorithmus:} Edmonds-Karp~\cite{edmonds1972, cormen2022} (BFS-Ford-Fulkerson), $O(VE^2)$.

% ============================================================
\section{Portfolio-Diversifikation: MST + Subgraph}
\label{sec:mst}
% ============================================================

\subsection{Korrelationsgraph und Pearson-Distanz}

Gegeben $n$ Assets mit Korrelationsmatrix $C \in [-1,1]^{n \times n}$.
Die \emph{Pearson-Distanz}~\cite{mantegna1999}:
\[
    d(i,j) = \sqrt{2(1 - \rho_{ij})} \in [0, 2]
\]

\textbf{Algorithmus:} Kruskal~\cite{kruskal1956, cormen2022} (Union-Find), $O(n^2 \log n)$.

\subsection{Portfolio-Strukturvergleich via Subgraph-Isomorphismus}

FyLab erlaubt den Vergleich zweier Portfolio-Korrelationsstrukturen:
\[
    \text{STRUKTURVERGLEICH} \leq_p \text{SUBGRAPH\_ISO}
\]

Die Korrelationsmatrix wird nach einem Schwellwert $\tau$ binarisiert:
\[
    A_{ij} = \begin{cases} 1 & |\rho_{ij}| > \tau \\ 0 & \text{sonst} \end{cases}
\]

Danach vergleicht \texttt{subgraph\_contains(A, B)} strukturell:
Ist die Abhängigkeitsstruktur von Portfolio $A$ in Portfolio $B$ enthalten?

\begin{lstlisting}[caption=Portfolio-Strukturvergleich]
from fylab.algorithms.mst import compare_portfolio_structures

result = compare_portfolio_structures(
    symbols_a, corr_a,  # Portfolio A
    symbols_b, corr_b,  # Portfolio B
    threshold=0.5,      # Schwelle fuer starke Korrelation
)
print(result.description)
# z.B. "A ist strukturell in B enthalten (B hat Superset der Abhaengigkeiten)"
\end{lstlisting}

\begin{satz}[MST als Diversifikationswerkzeug~\protect\cite{mantegna1999, markowitz1952}]
    Der MST des Korrelationsgraphen identifiziert Cluster ähnlicher Assets.
    Ein optimal diversifiziertes Portfolio wählt je einen Repräsentanten
    aus jedem MST-Cluster.
\end{satz}

% ============================================================
\section{Weitere Probleme: Steuer und Portfolio-Auswahl}
\label{sec:np_probleme}
% ============================================================

\subsection{Steueroptimierung: Subset-Sum}

\begin{definition}[Steueroptimierung]
    Gegeben Ausgaben $a_1, \ldots, a_n \in \mathbb{Z}_{>0}$ und ein
    absetzbares Budget $B$. Wähle eine Teilmenge $S \subseteq [n]$ mit
    $\sum_{i \in S} a_i \leq B$ und $\sum_{i \in S} a_i$ maximal.
\end{definition}

\begin{theorem}
    STEUER\_OPT ist NP-schwer (Reduktion aus SUBSET-SUM)~\cite{garey1979}.
\end{theorem}

\begin{beweis*}
    SUBSET-SUM $\leq_p$ STEUER\_OPT: Setze $B = T$ (Zielwert).
    SUBSET-SUM hat eine Lösung $\Leftrightarrow$ STEUER\_OPT liefert
    Summe $= B$.
\end{beweis*}

\textbf{Lösung:} Dynamische Programmierung in $O(nB)$ (pseudopolynomiell)~\cite{cormen2022}.

\subsection{Portfolioauswahl: 0-1-Rucksack}

\begin{reduktion}
    \textbf{PORTFOLIO\_SELECT $\leq_p$ KNAPSACK}

    Assets mit erwarteter Rendite $r_i$ (Nutzen), Kapital $p_i$ (Gewicht),
    Gesamtkapital $K$ (Kapazität). KNAPSACK mit FPTAS liefert~\cite{garey1979}
    $(1-\varepsilon)$-Approximation in $O(n \cdot K / \varepsilon)$.
\end{reduktion}

% ============================================================
\section{FyLab Modulhierarchie}
\label{sec:architektur}
% ============================================================

\begin{sidewaystable}
    \centering
    \begin{tabular}{@{}lll@{}}
        \toprule
        \textbf{Modul} & \textbf{Inhalt} & \textbf{Algorithmen} \\
        \midrule
        \texttt{core/}         & Konto, Portfolio, Cashflow, Währung
                               & -- \\
        \texttt{algorithms/}   & Subgraph, MST, MaxFlow, Arbitrage
                               & Subgraph Alg., Kruskal, Edmonds-Karp, Bellman-Ford \\
        \texttt{analysis/}     & Risikokennzahlen, Effizienzgrenze
                               & VaR, CVaR, Markowitz, Monte-Carlo \\
        \texttt{visualization/}& Charts, Graphen
                               & Matplotlib~\cite{hunter2007} \\
        \texttt{gui/}          & PyQt6-Fenster und Widgets
                               & -- \\
        \bottomrule
    \end{tabular}
    \caption{FyLab Modulhierarchie}
\end{sidewaystable}

\subsection{Klassifikation der GUI-Module analog PyLab}

\begin{sidewaystable}
    \centering
    \begin{tabular}{@{}lll@{}}
        \toprule
        \textbf{Widget} & \textbf{Funktion} & \textbf{Algorithmus} \\
        \midrule
        AccountWidget    & Buchungsverwaltung, Monatsübersicht & -- \\
        PortfolioWidget  & Depot, Allokation, Effizienzgrenze & Markowitz~\cite{markowitz1952}, Kruskal~\cite{kruskal1956} \\
        CashflowWidget   & Planung, Wasserfall-Chart & -- \\
        GraphWidget      & Betrug, Arbitrage, Schulden
                         & Subgraph Alg., Bellman-Ford, Greedy \\
        RiskWidget       & VaR, CVaR~\cite{rockafellar2000}, Monte-Carlo & Normalvert., Quantile~\cite{sharpe1964} \\
        \bottomrule
    \end{tabular}
    \caption{GUI-Widgets und zugehörige Algorithmen}
\end{sidewaystable}

% ============================================================
\section{Installation und Start}
\label{sec:install}
% ============================================================

\begin{lstlisting}[language=bash, caption=Installation und Start]
# FyLab installieren
pip install -e ".[dev]"

# Tests ausfuehren
pytest

# GUI starten
fylab-gui
\end{lstlisting}

% ============================================================
\section{Schluss}
\label{sec:schluss}
% ============================================================

FyLab demonstriert, dass die zentralen Strukturprobleme des Finanzwesens
direkt auf klassische Graphprobleme der Theoretischen Informatik
zurückgeführt werden können:

\begin{itemize}
    \item \textbf{NP-vollständige Probleme} (Subgraph-Isomorphismus,
    Portfolio-Auswahl, Steueroptimierung) werden durch Heuristiken,
    FPTAS oder Exponentialzeitalgorithmen mit polynomiellem Preprocessing
    behandelt.
    \item \textbf{Polynomial lösbare Probleme} (MST, MaxFlow, Bellman-Ford,
    Dijkstra) erlauben exakte Lösungen in vertretbarer Laufzeit.
    \item Der \textbf{Subgraph Algorithmus} (\texttt{git+gitlab.com/epp-group/subgraph})
    ist das algorithmische Herzstück von FyLab. Durch die zentrale Utility
    \texttt{subgraph\_contains\newline (R, G)} wird er in allen vier Algorithmusmodulen
    eingesetzt:
    \begin{enumerate}
        \item \texttt{subgraph\_finance}: Direkte Betrugsmuster-Erkennung, das heißt FRAUD\_\newline DETECT $\leq_p$ SUBGRAPH\_ISO
        \item \texttt{arbitrage}: Strukturelle Zyklus-Vorprüfung vor Bellman-Ford, das heißt CYCLE\_EXISTS $\leq_p$ SUBGRAPH\_ISO
        \item \texttt{max\_flow}: Schuldenstruktur-Mustererkennung, das heißt SCHULDEN\_\newline MUSTER $\leq_p$ SUBGRAPH\_ISO
        \item \texttt{mst}: Portfolio-Strukturvergleich zweier Korrelationsgraphen, das heißt STRUKTURVERGLEICH $\leq_p$ SUBGRAPH\_ISO
    \end{enumerate}
\end{itemize}

\subsection{Ausblick}

Das vorliegende FyLab-System bildet eine solide algorithmische Grundlage
für ein breites Spektrum weiterführender Erweiterungen.
Dieser Abschnitt gliedert mögliche Entwicklungsrichtungen entlang der
Dimensionen \emph{Algorithmen}, \emph{Architektur}, \emph{Daten},
\emph{Sicherheit} und \emph{Forschung}.

% ----------------------------------------------------------------
\subsection{Erweiterte Graphalgorithmen und theoretische Vertiefung}
\label{sub:ausblick_algo}
% ----------------------------------------------------------------

\paragraph{Dynamische Subgraph-Erkennung.}
Der aktuelle Subgraph Algorithmus~\cite{subgraph} operiert auf statischen
Adjazenzmatrizen. Finanztransaktionsgraphen sind jedoch inhärent zeitveränderlich:
Neue Buchungen entstehen, Kontosalden ändern sich, Betrugsmuster entwickeln sich
adaptiv. Ein nächster Schritt wäre die Integration eines
\emph{inkrementellen Subgraph-Matchings}~\cite{cordella2004}, das bei Graphänderungen
nur betroffene Teilstrukturen neu berechnet, statt den gesamten Graphen erneut
zu analysieren. Dies würde die amortisierte Laufzeit von $O(n^3)$ pro Update
auf $O(\Delta \cdot n^2)$ senken, wobei $\Delta$ die Anzahl geänderter Kanten
bezeichnet.

\paragraph{Hypergraph-Erweiterung für Mehrparteiengeschäfte.}
Klassische gerichtete Graphen modellieren bilaterale Transaktionen
(\emph{A überweist an B}). Viele reale Finanzvorgänge sind jedoch mehrparteienbasiert
(z.\,B.\ Syndikatkredite, DvP-Settlements, Escrow-Konstruktionen).
Eine \emph{Hypergraph-Repräsentation}, bei der Hyperkanten Mengen von
Teilnehmern verbinden, erlaubt eine präzisere Modellierung solcher Strukturen.
Die Subgraph-Isomorphismus-Suche auf Hypergraphen ist zwar noch NP-schwerer
als im Standardfall~\cite{garey1979}, jedoch erlauben polynomielle
Signaturen analoger Art eine effiziente Filterung in der Praxis.

\paragraph{Approximationsalgorithmen für Portfolio-Optimierung.}
Das Markowitz-Modell in \cite{markowitz1952} setzt vollständige Bekanntheit der
Kovarianzmatrix voraus. In der Praxis sind Korrelationen instabil und
müssen aus historischen Daten geschätzt werden. Hier bieten sich
\emph{robuste Optimierungsverfahren}~\cite{rockafellar2000} an, die explizit
Schätzunsicherheiten in die Zielfunktion einbeziehen: Min-Max-Portfolios,
ellipsoide Unsicherheitsmengen oder distributionell robuste Optimierung
(DRO) über Wasserstein-Kugeln.

\paragraph{Online-Algorithmen für Echtzeit-Arbitrage.}
Der aktuelle Bellman-Ford-Ansatz~\cite{bellman1958} berechnet negative Zyklen
im vollständigen Wechselkursgraphen. An Devisenmärkten ändern sich Kurse im
Millisekundenbereich. Ein \emph{Online-Algorithmus} müsste eingehende
Kurs-Updates in $O(\log n)$ verarbeiten und dabei negative Zyklen dynamisch
nachverfolgen – ein offenes Problem an der Schnittstelle von
Graphentheorie und Streaming-Computing.

\paragraph{Exakte Lösung des Knapsack-Problems via Branch-and-Bound.}
Das Portfolio-Auswahlproblem nutzt aktuell einen FPTAS-Ansatz~\cite{garey1979}.
Für Instanzen moderater Größe ($n \leq 500$ Assets) ist ein
\emph{Branch-and-Bound-Verfahren} mit Korrelationsschnitt als
obere Schranke oft exakter und schneller. Dies ist insbesondere für
Steueroptimierungsszenarien interessant, bei denen die Kapazität $B$
einem fixen Jahresfreibetrag entspricht und exakte Lösungen regulatorisch
gefordert werden.

% ----------------------------------------------------------------
\subsection{Maschinelles Lernen und KI-Integration}
\label{sub:ausblick_ml}
% ----------------------------------------------------------------

\paragraph{Graph Neural Networks für Betrugserkennung.}
Der regelbasierte Subgraph-Iso\-morphismus-Ansatz erkennt bekannte,
explizit modellierte Betrugsmuster (Wash-Trading, Smurfing, Layering).
\emph{Graph Neural Networks (GNNs)}~\cite{garey1979} können hingegen latente,
nicht vorab spezifizierte Anomalien lernen. Eine vielversprechende Architektur
wäre ein \emph{Graph Attention Network (GAT)}, das auf historisch gelabelten
Transaktionsgraphen trainiert wird und neue Betrugstypen ohne explizite
Muster-Definition detektiert. Der Subgraph Algorithmus könnte dabei als
effizienter \emph{Feature-Extraktor} dienen: erkannte Muster fließen als
binäre Merkmale in das GNN ein.

\paragraph{Reinforcement Learning für dynamische Asset-Allokation.}
Markowitz-basierte statische Portfoliooptimierung ignoriert die zeitliche
Dynamik von Märkten. Ein \emph{Deep Reinforcement Learning (DRL)}-Agent
könnte täglich Portfoliogewichte anpassen, basierend auf dem aktuellen
MST-Clusteringzustand des Korrelationsgraphen. Der MST dient als
\emph{Zustandsrepräsentation}: sein Topologiewechsel (z.\,B.\ Verlust einer
Hub-Kante) signalisiert Regimewechsel am Markt und löst Rebalancing-Aktionen
aus.

\paragraph{Anomalieerkennung in Cashflow-Zeitreihen.}
Isolierte Kassenabweichungen, unerwartete Frequenzänderungen und saisonale
Drift sind schwer durch statische Regeln zu erfassen. \emph{Autoregressive
Transformermodelle} (analogie zu GPT für Zeitreihen) könnten auf dem
Cashflow-Planungsmodul trainiert werden, um Abweichungen von der erwarteten
monatlichen Dynamik statistisch zu quantifizieren und dem Nutzer
frühzeitig anzuzeigen.

% ----------------------------------------------------------------
\subsection{Persistenz, Import und Datenintegration}
\label{sub:ausblick_daten}
% ----------------------------------------------------------------

\paragraph{Persistente Datenhaltung.}
Version~1.0 hält alle Daten im Arbeitsspeicher. Ein produktives System
benötigt eine eingebettete Datenbank (SQLite via \texttt{SQLAlchemy} oder
DuckDB) mit klar definiertem Schema für Konten, Transaktionen, Positionen
und Cash\-flow-Pläne. Das Datenbankschema würde dabei die algebraische
Struktur der Graph-Adjazenz\-¸matrizen widerspiegeln: Knoten als Entitäten-
tabellen, Kanten als Relationsatbellen mit FK-Referenzen.

\paragraph{CSV/JSON/MT940-Import.}
Reale Kontoauszüge liegen in Formaten wie \texttt{MT940},
\texttt{OFX/QFX} (Quicken) oder bankenspezifischen CSV-Exporten vor.
Eine robuste Import-Pipeline mit automatischer Kategorisierung
(via Schlüsselwort-Heuristik oder kleinem Klassifikationsmodell)
und Duplikaterkennung würde FyLab für den produktiven Einsatz
qualifizieren. Dies ist als Milestone 1.1 im Quellcode markiert.

\paragraph{Live-Marktdaten via REST-APIs.}
Aktuelle Aktienkurse, Wechselkurse und Zinssätze könnten über
Standard-APIs (Yahoo Finance, Alpha Vantage, ECB Data Portal) abgerufen
und direkt in die Algorithmusmodule eingespeist werden. Der Arbitrage-Detektor
würde so zu einem \emph{Echtzeit-Alert-System}: bei Überschreitung einer
vordefinierten Arbitrage-Schwelle ($\prod_k r_k > 1 + \varepsilon$) generiert
FyLab automatisch eine Handlungsempfehlung.

\paragraph{Blockchain-Transaktionsgraphen.}
Die Graph-Infrastruktur von FyLab ist nicht auf klassische Bankbuchungen
beschränkt. Öffentliche Blockchain-Ledger (Ethereum, Bitcoin) exponieren
Transaktionsgraphen, auf denen die Betrugsmuster-Erkennung via Subgraph-
Isomorphismus unmittelbar anwendbar ist. Mixer-Muster, Tornado-Cash-ähnliche
Strukturen und Chain-Hopping könnten als spezialisierte \texttt{FraudPattern}-
Einträge modelliert und automatisch erkannt werden.

% ----------------------------------------------------------------
\subsection{Architektur und Skalierbarkeit}
\label{sub:ausblick_arch}
% ----------------------------------------------------------------

\paragraph{Asynchrone Berechnung im GUI.}
Langläufige Operationen (Monte-Carlo mit $10^6$ Pfaden, MST auf
500-Asset-Portfolios, Effizienzgrenzberechnung mit $10^4$ Punkten)
blockieren aktuell den Qt-Event-Loop. Die Umstellung auf
\texttt{QThread} oder \texttt{asyncio}-kompatible Qt-Coroutinen
(über \texttt{qasync}) würde responsives UI-Verhalten auch bei rechenintensiven
Analysen gewährleisten ohne strukturellen Umbau der Algorithmusmodule.

\paragraph{Plugin-Architektur für Algorithmen.}
Das aktuelle System registriert Algorithmen statisch im \texttt{algorithms/}-Pakket.
Eine \emph{Plugin-API} (ähnlich dem pytest-Plugin-System) würde es erlauben,
neue Algorithmusmodule zur Laufzeit zu laden: Drittentwickler könnten eigene
Betrugsmuster, Risikokennzahlen oder Optimierungsverfahren als
\texttt{fylab-plug\newline in-*}-Pakete publizieren, die automatisch in die Menüs und
Widgets des Hauptfensters integriert werden.

\paragraph{Parallele Subgraph-Suche.}
Die $O(n^3)$-Hauptschleife des Subgraph Algorithmus~ \cite{subgraph}
ist datenparallel: Signaturen verschiedener Knotenpaare sind unabhängig
berechenbar. Mit \texttt{numpy}-Vektorisierung (bereits teilweise realisiert)
oder einer CUDA-Kernel-Implemen\-tierung (via \texttt{cupy}) ließe sich die
Laufzeit auf Portfolios mit $n > 1000$ Knoten auf Millisekunden reduzieren.

\paragraph{REST-API und Microservice-Architektur.}
Eine \texttt{FastAPI}-basierte REST-Schicht über den Algorithmusmodulen würde
FyLab als Backend-Service für Web-Frontends, Mobile Apps oder
institutionelle Integrationspipelines positionieren. Jedes Modul
(\texttt{/api/v1/\newline arbitrage}, \texttt{/api/v1/fraud}, \texttt{/api/v1/mst},
\texttt{/api/v1/portfolio}) exponiert seine Kernfunktionalität als
zustandslosen HTTP-Endpunkt, der in Kubernetes horizontal skalierbar ist.

% ----------------------------------------------------------------
\subsection{Regulatorik, Compliance und Sicherheit}
\label{sub:ausblick_regulatorik}
% ----------------------------------------------------------------

\paragraph{AML/KYC-Compliance-Engine.}
Die Betrugsmuster-Erkennung von FyLab tangiert den regulatorischen Bereich
\emph{Anti-Money Laundering (AML)} und \emph{Know Your Customer (KYC)}.
Eine dedizierte Compliance-Erweiterung würde FATF-Rekommendationen
\cite{garey1979}\footnote{Die formale Komplexitätstheorie ist hier analog:
AML-Regelwerke definieren verbotene Transaktionsmuster, deren Erkennung
im Transaktionsgraphen ein Subgraph-ISO-Problem darstellt.}
als formale Muster in \texttt{FraudPattern}-Typen übersetzen und
einen revisionssicheren Audit-Trail aller Erkennungsergebnisse führen.

\paragraph{Differenzielle Privatheit für Portfolio-Analysen.}
Portfolioanalysen offenbaren sensible Vermögensdaten. Durch Integration
von \emph{differenzieller Privatheit} ($\varepsilon$-DP) in den
Monte-Carlo-Simulator und die Effizienzgrenzenberechnung könnten
Ergebnisse mit garantiertem Datenschutz an Dritte (z.\,B.\ Steuerberater,
Banken) weitergegeben werden, ohne individuelle Positionen zu exponieren.

\paragraph{Kryptographische Verifikation.}
Transaktionsdaten in der eingebetteten Datenbank könnten mittels
kryptographischer Hashketten (Merkel-Baum über sortierten Buchungen) gegen
nachträgliche Manipulation gesichert werden – eine leichtgewichtige
Auditierbarkeit abseits vollständiger Blockchain-Lösungen.

% ----------------------------------------------------------------
\subsection{Wissenschaftlicher Forschungsausblick}
\label{sub:ausblick_forschung}
% ----------------------------------------------------------------

Die in dieser Arbeit etablierten polynomiellen Reduktionen öffnen mehrere offene Forschungsfragen:

\begin{enumerate}
    \item \textbf{Schärfe der Subgraph-Isomorphismus-Schranke.}\\
    Gibt es finanzmathematisch relevante Teilklassen des
    Subgraph-Isomorphismus-Pro\-blems (z.\,B.\ planare Finanzgraphen,
    Baumstruktur-Portfolios, DAG-Schuldennetz\-werke), für die polynomielle
    Algorithmen mit garantierter Korrektheit existieren? Die Theorie
    liefert hier: Subgraph-Isomorphismus auf Bäumen ist in
    $O(n^{1.5})$ lösbar~\cite{cormen2022}.

    \item \textbf{Durchschnittliche Laufzeit auf realen Finanzdaten.}\\
    Der Subgraph Algorithmus~\cite{subgraph} hat Worst-Case
    $O(n^3)$, verhält sich empirisch jedoch nahe $O(n^2)$ auf
    typischen Finanzgraphen (spärliche Adjazenzmatrizen, wenige
    Zyklen). Eine formale Average-Case-Analyse über Random-Graph-Modelle
    (Erdős-Rényi, Barabási–Albert Skalenfreiheit) wäre publikationswürdig.

    \item \textbf{Untere Komplexitätsschranken für Betrugserkennungs-Varianten.}\\
    Welche spezifischen Betrugsmuster-Erkennungsprobleme sind in $P$?
    Wash-Trading auf bipartiten Graphen? Smurfing auf Sternstrukturen?
    Eine vollständige Klassifikation der Tracklabeled-Subgraph-
    Isomorphismus-Varianten nach ihrer Komplexität fehlt noch in der Literatur.

    \item \textbf{Algebraische Charakterisierung der Pearson-Distanz-MSTs.}\\
    Mantegna~\cite{mantegna1999} zeigte empirisch, dass Pearson-MSTs
    stabile Branching-Strukturen aufweisen. Unklar ist, welche
    algebraischen Eigenschaften des Korrelationstensors diese Stabilität
    garantieren – eine Verbindung zur Spektraltheorie positiv definiter
    Matrizen und zum Fiedler-Vektor des Laplacians wäre mathematisch
    fruchtbar.

    \item \textbf{Verbindung zu Algorithmic Game Theory.}\\
    Schuldenausgleich via Max-Flow ist ein spezieller Fall von
    multi-commodity-flow. Wenn Agenten strategisch falsche Schulden
    deklarieren, entsteht ein unvollständiges-Informations-Spiel.
    Die Nash-Gleichgewichtsanalyse dieses Spiels und die Frage,
    ob wahrheitsgemäße Deklaration eine dominante Strategie ist
    (Mechanism Design), ist ein offenes Problem.
\end{enumerate}

\begin{table}[h]
\centering
\caption{Zusammenfassung der Entwicklungs- und Forschungsrichtungen}
\label{tab:ausblick}
\small
\begin{tabular}{@{}lll@{}}
    \toprule
    \textbf{Bereich} & \textbf{Erweiterung} & \textbf{Zeithorizont} \\
    \midrule
    Algorithmen   & Inkrementelles Subgraph-Matching           & kurzfristig \\
    Algorithmen   & GNN-gestützte Betrugserkennung             & mittelfristig \\
    Algorithmen   & DRL für dynamische Asset-Allokation        & mittelfristig \\
    Daten         & MT940/CSV-Import-Pipeline                  & kurzfristig \\
    Daten         & Live-API-Anbindung (ECB, Alpha Vantage)    & kurzfristig \\
    Daten         & Blockchain-Transaktionsgraph               & mittelfristig \\
    Architektur   & Qt-Threads für Async-Berechnung            & kurzfristig \\
    Architektur   & REST-API / Microservice                    & mittelfristig \\
    Architektur   & Plugin-System für Drittmodule              & langfristig \\
    Sicherheit    & AML/KYC-Compliance-Engine                  & mittelfristig \\
    Sicherheit    & Differenzielle Privatheit                  & langfristig \\
    Forschung     & Average-Case-Analyse Subgraph-ISO          & mittelfristig \\
    Forschung     & Mechanism Design Schuldenausgleich         & langfristig \\
    \bottomrule
\end{tabular}
\end{table}

\begin{thebibliography}{99}

\bibitem{subgraph}
Stephan Epp.
\textit{Der Subgraph Algorithmus}, 7.\ Januar 2026.
Verfügbar unter: \url{https://gitlab.com/epp-group/subgraph}

\bibitem{ullmann1976}
J.~R.~Ullmann.
\newblock An algorithm for subgraph isomorphism.
\newblock \textit{Journal of the ACM}, 23(1):31--42, 1976.

\bibitem{cordella2004}
L.~P.~Cordella, P.~Foggia, C.~Sansone, M.~Vento.
\newblock A (Sub)Graph Isomorphism Algorithm for Matching Large Graphs.
\newblock \textit{IEEE Transactions on Pattern Analysis and Machine
  Intelligence}, 26(10):1367--1372, 2004.

\bibitem{cook1972}
S.~A.~Cook.
\newblock The complexity of theorem-proving procedures.
\newblock In \textit{Proceedings of the 3rd Annual ACM Symposium on Theory
  of Computing (STOC)}, S.~151--158. ACM, New York, 1972.

\bibitem{garey1979}
M.~R.~Garey, D.~S.~Johnson.
\newblock \textit{Computers and Intractability: A Guide to the Theory of
  NP-Completeness}.
\newblock W.~H.~Freeman, San Francisco, 1979.

\bibitem{cormen2022}
T.~H.~Cormen, C.~E.~Leiserson, R.~L.~Rivest, C.~Stein.
\newblock \textit{Introduction to Algorithms}, 4.~Auflage.
\newblock MIT Press, Cambridge, MA, 2022.

\bibitem{bellman1958}
R.~Bellman.
\newblock On a routing problem.
\newblock \textit{Quarterly of Applied Mathematics}, 16(1):87--90, 1958.

\bibitem{ford1956}
L.~R.~Ford, D.~R.~Fulkerson.
\newblock Maximal flow through a network.
\newblock \textit{Canadian Journal of Mathematics}, 8:399--404, 1956.

\bibitem{edmonds1972}
J.~Edmonds, R.~M.~Karp.
\newblock Theoretical improvements in algorithmic efficiency for network
  flow problems.
\newblock \textit{Journal of the ACM}, 19(2):248--264, 1972.

\bibitem{kruskal1956}
J.~B.~Kruskal.
\newblock On the shortest spanning subtree of a graph and the traveling
  salesman problem.
\newblock \textit{Proceedings of the American Mathematical Society},
  7(1):48--50, 1956.

\bibitem{dijkstra1959}
E.~W.~Dijkstra.
\newblock A note on two problems in connexion with graphs.
\newblock \textit{Numerische Mathematik}, 1(1):269--271, 1959.

\bibitem{mantegna1999}
R.~N.~Mantegna.
\newblock Hierarchical structure in financial markets.
\newblock \textit{The European Physical Journal B}, 11(1):193--197, 1999.

\bibitem{markowitz1952}
H.~Markowitz.
\newblock Portfolio selection.
\newblock \textit{The Journal of Finance}, 7(1):77--91, 1952.

\bibitem{sharpe1964}
W.~F.~Sharpe.
\newblock Capital asset prices: A theory of market equilibrium under
  conditions of risk.
\newblock \textit{The Journal of Finance}, 19(3):425--442, 1964.

\bibitem{rockafellar2000}
R.~T.~Rockafellar, S.~Uryasev.
\newblock Optimization of conditional value-at-risk.
\newblock \textit{Journal of Risk}, 2(3):21--41, 2000.

\bibitem{artzner1999}
P.~Artzner, F.~Delbaen, J.-M.~Eber, D.~Heath.
\newblock Coherent measures of risk.
\newblock \textit{Mathematical Finance}, 9(3):203--228, 1999.

\bibitem{glasserman2004}
P.~Glasserman.
\newblock \textit{Monte Carlo Methods in Financial Engineering}.
\newblock Springer, New York, 2004.

\bibitem{harris2020}
C.~R.~Harris, K.~J.~Millman, S.~J.~van der Walt \textit{et al.}
\newblock Array programming with NumPy.
\newblock \textit{Nature}, 585:357--362, 2020.

\bibitem{virtanen2020}
P.~Virtanen, R.~Gommers, T.~E.~Oliphant \textit{et al.}
\newblock SciPy 1.0: Fundamental algorithms for scientific computing
  in Python.
\newblock \textit{Nature Methods}, 17:261--272, 2020.

\bibitem{hunter2007}
J.~D.~Hunter.
\newblock Matplotlib: A 2D graphics environment.
\newblock \textit{Computing in Science \& Engineering}, 9(3):90--95, 2007.

\bibitem{riverbank2023}
Riverbank Computing Limited.
\newblock \textit{PyQt6 Reference Guide}, 2023.
\newblock \url{https://www.riverbankcomputing.com/static/Docs/PyQt6/}.
  Letzter Zugriff: März 2026.

\end{thebibliography}

\end{document}
